% !TeX root = ../../Book5.tex
% 纲要标题：### 8.1 半单性的失效
% 纲要位置：工作记录/计划纲要.md，第 3983 行。

\section{半单性的失效}
\label{b5:sec:0801}


% 纲要范围：
%
% * 特征 p 整除群阶
% * 群代数的根及不可分解模
% * 循环 p 群的完整基本例子
% * 中心 p 元在简单模上的作用；一般群代数的根包含于增广理想，p 群时相等
%

当底域的特征整除群阶时，平均投影中的除法不再有意义。表示仍然可以有合成列，但这些简单商之间的扩张未必分裂；因此研究简单模还不够，还须研究不可分解模及其根层。本章以循环 \(p\) 群为反复计算的例子，始终区分直和分解与合成因子。

\subsection{\texorpdfstring{特征 \(p\) 整除群阶的情形}{特征 p 整除群阶的情形}}
\label{b5:subsec:0801:01}

\begin{definition}\label{b5:def:modular-representation}
设 \(G\) 为有限群，\(k\) 为特征 \(p>0\) 的域。\(G\) 的有限维 \(k\) 表示称为特征 \(p\) 的\term[mo-biaoshi]{模表示}[modular representation]。本章重点讨论 \(p\,\bigm\vert\,|G|\) 的情形；若 \(p\mathrel{\vcenter{\hbox{\scalebox{1.25}{\(\nmid\)}}}}|G|\)，\cref{b5:thm:maschke}仍保证完全可约。
\end{definition}
由\cref{b5:prop:maschke-converse}，当 \(p\,\bigm\vert\,|G|\) 时，增广映射
\[
\varepsilon:kG\longrightarrow k,\qquad \sum_g a_gg\longmapsto\sum_g a_g
\]
没有 \(kG\) 线性截面。具体地，一个不动元必须为 \(c\sum_g g\)，它的增广为 \(c|G|=0\)，所以不能提升目标中的一。这既给出不完全可约的正则模，也给出取不动空间不保持满射的例子。
\ProseParagraphBreak
最小的矩阵例子是 \(C_p=\langle s\rangle\) 在 \(k^2\) 上作用为
\[
s\longmapsto\begin{pmatrix}1&1\\0&1\end{pmatrix}.
\]
第一坐标轴和商都为平凡表示，但没有不变补。这个表示与平凡二维表示有相同的两个简单合成因子，却不同构：前者的生成元不是恒等。后面定义的 Brauer 特征标会记录合成因子的重数，不会额外记录这条扩张是否分裂。

\subsection{群代数的根与不可分解模}
\label{b5:subsec:0801:02}

% 跨卷引用目标：b2:thm:module-radical-jm（第二卷定理15.2.2）。
设 \(A\) 为有限维含幺 \(k\) 代数。其 Jacobson 根可写为全部简单左模的零化子的交，记为 \(J=J(A)\)。若 \(M\) 为有限维左模，全部极大子模的交称为其\term[mo-gen]{根}[radical of a module]。它等于 \(JM\)，如第二卷定理15.2.2所证：每个极大商都是被 \(J\) 零化的简单模，故 \(JM\) 包含在每个极大子模中；反过来，\(M/JM\) 半单，任一非零商类都能在某个简单分量上得到非零投影，该投影的核给出一个不含原元素的极大子模。商 \(M/JM\) 称为\term[mo-ding]{顶}[top of a module]。根层 \(J^iM/J^{i+1}M\) 记录从顶逐步向下的简单成分。它们通常不能在原模中同时选成直和项；即使全部根层逐项同构，原模也未必同构。\subsecref{b5:subsec:0804:03}将用 \(C_2\times C_2\) 的一参数模族计算这种现象：各层相同，层间的作用仍可不同。
\symidx{\(J(A),JM,M/JM\)：代数的 Jacobson 根与模的根、顶}
\begin{proposition}\label{b5:prop:finite-radical-layers}
设 \(k\) 为域，\(A\) 为有限维含幺 \(k\) 代数，\(J=J(A)\)。则 \(J\) 幂零，\(A/J\) 半单。对任意有限维左 \(A\) 模 \(M\)，全部根层 \(J^iM/J^{i+1}M\) 都是半单模，且 \(JM=M\) 蕴涵 \(M=0\)。
\end{proposition}
\begin{proof}
% 跨卷目标：b2:prop:finite-algebra-semisimple-quotient（第二卷命题15.1.3）。
取左正则模 \(A\) 的有限合成列 \(0=A_0\subset A_1\subset\cdots\subset A_\ell=A\)。\(J\) 零化每个简单商 \(A_i/A_{i-1}\)，所以 \(JA_i\subseteq A_{i-1}\)。逐次相乘便有 \(J^rA\subseteq A_{\ell-r}\) 对 \(0\le r\le\ell\) 成立，特别地 \(J^\ell A\subseteq A_0=0\)，即 \(J^\ell=0\)。有限维代数的半单商结构由第二卷命题15.1.3及该节建立的半单商分解给出：\(A/J\) 为有限个除环上矩阵代数的直积。每个根层被 \(J\) 零化，是 \(A/J\) 模，因而半单。最后若 \(JM=M\)，则 \(M=J^\ell M=0\)。
\end{proof}
上述最后一步是这里需要的 Nakayama 引理\footnote{中山正（假名：なかやま ただし，罗马音：Tadashi Nakayama，1912-1964），日本数学家，研究环论和群表示论。}。它也说明为什么顶不能为零：任何非零有限维模都有非零的顶。有限维模按不可分解项的直和分解及其唯一性可用第二卷定理12.4.2；该定理的有限长度条件在此由有限维性满足。不可分解模的各根层则可能包含多种简单模。

\ProseParagraphBreak
对任意有限群 \(G\) 和域 \(k\)，平凡一维模 \(k\) 都是简单左 \(kG\)-模，而 \(a=\sum_g a_gg\) 在它上面的作用是乘以 \(\varepsilon(a)=\sum_g a_g\)。因此
\[
J(kG)=\bigcap_{S\;\text{简单}}\operatorname{Ann}_{kG}(S)
\subseteq\operatorname{Ann}_{kG}(k)=\ker\varepsilon.
\]
增广理想只须零化平凡模，根却须零化全部简单模，所以这条包含关系未必是等号。下面的有限 \(p\)-群情形中，唯一的简单类型就是平凡模，因而两个理想相等。
\begin{proposition}\label{b5:prop:p-group-local}
设 \(p\) 为素数，\(P\) 为有限 \(p\)-群，\(k\) 为特征 \(p\) 的域，\(\varepsilon:kP\to k\) 为增广映射 \(\sum_g a_gg\mapsto\sum_g a_g\)。唯一的简单左 \(kP\)-模是平凡一维模 \(k\)，且
\[
J(kP)=\ker\varepsilon,\qquad kP/J(kP)\cong k.
\]
因此 \(kP\) 是局部代数。
\end{proposition}
\begin{proof}
按 \(|P|\) 归纳。\(P=1\) 时，简单 \(k\)-向量空间只有一维模。若 \(P\ne1\)，第一卷定理4.2.5给出 \(p\,\bigm\vert\,|Z(P)|\)：类方程中各非中心类的大小都被 \(p\) 整除。取中心中一个阶为 \(p^a>1\) 的元，再取其 \(p^{a-1}\) 次幂，便得阶恰为 \(p\) 的元 \(z\)。在简单模 \(S\) 上令 \(F(v)=(z-1)v\)。对 \(g\in P\)、\(v\in S\)，中心性给出
\[
gF(v)=g(zv-v)=z(gv)-gv=F(gv).
\]
因此 \(F\) 属于模自同态环 \(\operatorname{End}_{kP}(S)\)，而不只是一个线性算子。又有 \(F^p=0\)，因为特征为 \(p\) 且 \((z-1)^p=z^p-1=0\)。由\cref{b5:lem:schur}，简单模的每个非零模自同态都可逆；幂零算子不能可逆，所以 \(F=0\)。于是 \(z\) 平凡作用，\(S\) 的作用通过较小的 \(P/\langle z\rangle\) 给出，归纳使它平凡。平凡作用下简单性又迫使维数为一。根是这个唯一简单模的零化子，恰为增广理想。模掉根是域，所以非单位元恰为根，代数局部。
\end{proof}
中心性在上述证明中不可少。取特征二的域 \(k\)，\cref{b5:exm:s3-two-blocks}将构造 \(S_3=\langle r,s\rangle\) 的二维简单模，其作用矩阵为
\[
R=\rho(r)=\begin{pmatrix}0&1\\1&1\end{pmatrix},\qquad
S=\rho(s)=\begin{pmatrix}0&1\\1&0\end{pmatrix}.
\]
换位 \(s\) 阶为二，且 \((S-I)^2=0\)，但 \(S-I\ne0\)。在这个模上直接计算得到
\[
(S-I)R-R(S-I)=\begin{pmatrix}0&1\\1&0\end{pmatrix}\ne0,
\]
所以 \(S-I\) 不属于模自同态环，不能对它应用 Schur 引理。非中心群元素的作用可能恰好与其他作用交换，但其阶为 \(p\) 所给出的幂零性本身不保证这一点。

\ProseParagraphBreak
例如一个非循环 \(p\) 群仍然只有一个简单模，但其不可分解模远不止一种。简单类型的数目很少，并不意味着扩张和直和分解的分类也很短。

\subsection{\texorpdfstring{循环 \(p\) 群的基本例子}{循环 p 群的基本例子}}
\label{b5:subsec:0801:03}

\begin{theorem}\label{b5:thm:cyclic-p-modules}
设 \(p\) 为素数，\(k\) 为特征 \(p\) 的域，\(r\ge1\) 为整数，\(G=C_{p^r}=\langle s\rangle\)，并记 \(m=p^r\)。则
\[
kG\cong k[t]/(t^m),\qquad t\longmapsto s-1.
\]
全部有限维不可分解左 \(kG\)-模的同构类代表恰为
\[
M_d=k[t]/(t^d),\qquad 1\le d\le m.
\]
每个有限维左 \(kG\)-模唯一分解成这些模的直和，唯一性指各项同构类的重数唯一。记 \(J=J(kG)\)，则对 \(i\ge0\) 及 \(1\le a,b,d\le m\) 有
\[
J^iM_d=t^iM_d,\qquad
\operatorname{soc}(M_d)=kt^{d-1},\qquad
\dim_k\operatorname{Hom}_{kG}(M_a,M_b)=\min(a,b).
\]
其中\term[jizuo]{基座}[socle] \(\operatorname{soc}(M_d)\) 是全部简单子模之和。
\end{theorem}
\symidx{\(\operatorname{soc}(M_d)\)：模的基座，即全部简单子模之和}
\begin{proof}
特征 \(p\) 中重复使用 \((x-1)^p=x^p-1\)，得到 \((s-1)^m=0\)。由 \(s=1+t\) 可知所给代数同态满射，两侧维数都是 \(m\)，故同构。模结构因此由满足 \(T^m=0\) 的线性算子 \(T\) 决定。幂零算子的 Jordan 链分解，或第二卷定理4.3.2在 \(k[t]\) 上的应用，给出模为 \(k[t]/(t^d)\) 的直和，且 \(d\le m\)。这里多项式只有不可约因子 \(t\)，所以不要求 \(k\) 代数闭。
\(M_d\) 的自同态由一的像决定，环为 \(k[t]/(t^d)\)。这个环的常数项非零元素可用有限几何级数求逆，常数项零元素幂零，故它局部，没有非平凡幂等元。因此 \(M_d\) 不可分解；不同 \(d\) 的维数不同。任一分解中，大小至少为 \(j\) 的块数等于
\[
\dim\ker T^j-\dim\ker T^{j-1},
\]
所以块的重数唯一。根层由 \(J=(t)\) 直接给出；简单子模被 \(t\) 零化，而 \(\ker(t:M_d\to M_d)=kt^{d-1}\)，得到基座。最后，\(M_a\to M_b\) 的同态由一个满足 \(t^a v=0\) 的 \(v\in M_b\) 决定，这个核的维数为 \(\min(a,b)\)。
\end{proof}
例如在本定理的设定中，另假定 \(m=p^r\ge4\)，令 \(A=k[t]/(t^m)\)。此时 \(M_2\) 与 \(M_4\) 都是所列的 \(A\)-模。若 \(f:M_2\to M_4\) 是模同态，则必须有 \(t^2f(1)=0\)，所以
\[
f(1)\in\ker(t^2:M_4\to M_4)=kt^2\oplus kt^3.
\]
反过来，\(g:M_4\to M_2\) 所需的条件 \(t^4g(1)=0\) 对 \(M_2\) 的每个元素都成立，故
\[
g(1)\in M_2=k1\oplus kt
\]
可以任意选取。两处 Hom 空间的维数都是二，但前一方向的像必须落在较长链的末端 \(t^2M_4\)，反向则允许生成元映到短链中的任意元素。例如 \(f(1)=t^2\) 给出单射，而 \(g(1)=1\) 给出满射；维数公式的对称性没有抹去这两种方向的区别。

\ProseParagraphBreak
这些模的全部子模也成一条链。因为 \(M_d\) 是 \(k[t]\) 的循环商模，子模对应包含 \((t^d)\) 的理想，恰为 \((t^j)\)、\(0\le j\le d\)。例如 \(M_3\) 的根层各为一个平凡模，却不是三个平凡模的直和；其不动空间只有一维，后者有三维。
